\emph{The conclusion.} ``$\ell\nmid h_{3,n}/h_{3,n-1}$'' is read as $v_\ell(h_{3,n})=v_\ell(h_{3,n-1})$, equivalently $\ker(N)_\ell=0$ for the norm $N$ of $B_{3,n}/B_{3,n-1}$ (Lemma~\ref{lem:oldnew} with $p=3$, $\ell\ne3$). \emph{The prime.} Since $\ell$ is odd and $\ell\ne3$, the two places where the oddness of $\ell$ enters are available: Lemma~\ref{lem:normone}(iv) (the transfer of relative norm $1$ to the $\ell$-th root, where $\operatorname{Nr}(\varepsilon)^{\ell}=1$ with $\operatorname{Nr}(\varepsilon)$ real forces $\operatorname{Nr}(\varepsilon)=1$, \path{WeberP3Rel.eq_one_of_odd_pow_eq_one}) and the hypothesis of Theorem~\ref{thm:SH}, which is stated for an odd prime. Nothing below is claimed for $\ell=2$.

\emph{The data.} Let $p=3$, $N=3^n$, $q=3^{n+1}$, $\zeta=e^{2\pi i/q}$, and let $\sigma$ be the generator $\zeta\mapsto\zeta^{4}$ of $\operatorname{Gal}(\Q(\zeta)/\Q(\zeta^{3^n}))\cong\Z/3^n$, which restricts to a generator of $G=\operatorname{Gal}(B_{3,n}/\Q)$. The Horie unit is $\eta_n=\sin(2(1+3^n)\pi/q)/\sin(2\pi/q)$ (\cite{MO13} \S1; at $p=3$ the product defining $\eta_n$ has one factor), a unit of $B_{3,n}$ of relative norm $1$ (Lemma~\ref{lem:normone}). For $\alpha=\sum_{i<c_n}a_i\zeta_{3^n}^i\in\Z[\zeta_{3^n}]$ ($c_n=2\cdot3^{n-1}$, $\zeta_{3^n}=\zeta^3$) put $\alpha_\sigma=\sum_{i<c_n}a_i\sigma^i$ and $u_\alpha=\eta_n^{\alpha_\sigma}$; for $\alpha\in\Z[\zeta_{3^r}]\subseteq\Z[\zeta_{3^n}]$ (via $\zeta_{3^r}=\zeta_{3^n}^{3^{n-r}}$) the exponents of $\alpha_\sigma$ are the multiples $3^{n-r}j$, $j<c$. The map $\alpha\mapsto H(u_\alpha)$ is additive in the coordinates $a_j$, and its image on $\Z[\zeta_{3^r}]$ is the lattice spanned by $H(\sigma^{3^{n-r}j}\eta_n)$, $j<c$, whose covolume $D_r^{(n)}$ is positive by Theorem~\ref{thm:rank3}: this is (Rank) of Theorem~\ref{thm:SH}, for every $n$ and $r$.

\emph{Saturation (the input of \cite{MO13} Lemma~1.3).} Since $r\le s$, the field $F=\Q(\zeta_{3^r})$ contains the decomposition field of $\ell$ in $\Q(\zeta_{3^n})$ (for $f=1$ the decomposition field is $\Q(\zeta_{3^{\min(n,s)}})$ itself; for $f=2$ it is its maximal real subfield). If $\ell\mid h_{3,n}/h_{3,n-1}$, \cite{MO13} Lemma~1.3 (attributed there to \cite{Ho05b}, which relays \cite{Ho02}; the general form with proof is \cite{Ho05a} Prop.~1) gives a prime $\mathcal L$ of $F$ above $\ell$ such that $u_\alpha=\eta_n^{\alpha_\sigma}$ is an $\ell$-th power in $E_{3,n}$ for every $\alpha\in\ell\mathcal L^{-1}$. The residue degree of $\mathcal L$ is $f$, so $\ell\mathcal L^{-1}$ has index $\ell^{c-f}$ in $\Z[\zeta_{3^r}]$: this is (Index) with $d=f$. Let $U=U_{3,n}=\{\varepsilon\in E_{3,n}:\operatorname{Nr}_{B_{3,n}/B_{3,n-1}}\varepsilon=1\}$. (Car) for $\alpha\in\ell\mathcal L^{-1}\setminus\ell\Z[\zeta_{3^r}]$: the $\ell$-th root $\varepsilon$ has $\ell H(\varepsilon)=H(u_\alpha)$ (\path{WeberSH.sat_root_log}) and lies in $U_{3,n}$ by Lemma~\ref{lem:normone}(iv). (Car) for $\alpha=\ell\beta$: $u_\beta$ itself carries $H(u_\alpha)=\ell H(u_\beta)$ (\path{WeberSH.base_branch}), and $u_\beta\in U_{3,n}$ because $\operatorname{Nr}(\eta_n^{\beta_\sigma})=\operatorname{Nr}(\eta_n)^{\beta_\sigma}=1$ (Lemma~\ref{lem:normone}(iii)).

\emph{The floor.} Let $\varepsilon\in U_{3,n}$ with $H(\varepsilon)\ne0$. By definition of $U_{3,n}$, $\operatorname{Nr}_{B_{3,n}/B_{3,n-1}}(\varepsilon)=1$; by Lemma~\ref{lem:normone}(v) (a unit of $E_{3,n-1}$ with relative norm $1$ equals $1$, and $H(1)=0$), $\varepsilon\notin E_{3,n-1}$. Hence $\varepsilon\in E_{3,n}\setminus E_{3,n-1}$ with relative norm $1$, and Lemma~\ref{lem:mo25-z3} gives $\hgt(\varepsilon)\ge L^{\mathrm{rel}}_{3,n}$ directly in the Euclidean height (the height of \cite{MO16} Def.~2.2 is the $\ell^2$-norm of the log vector; no Cauchy--Schwarz step is needed). This is (Floor) for $U=U_{3,n}$ with $L_0=L^{\mathrm{rel}}_{3,n}$.

\emph{Conclusion.} Theorem~\ref{thm:SH} (\path{WeberSH.theoremSH_contra}; \path{WeberP3.theoremP3_core} is the same statement with the $p=3$ names) with $K=B_{3,n}$, $m=c$, the lattice $\ell\mathcal L^{-1}$, $L_0=L^{\mathrm{rel}}_{3,n}$ and Blichfeldt's theorem (\cite{Bl14} Thm.~II as \cite{MO16} Thm.~2.7) on $\Lambda=\ell^{-1}H(u_{\ell\mathcal L^{-1}})$, of rank $c$ and covolume $D_r^{(n)}/\ell^{f}$ in its span, yields $\ell^{f}\le(2/\pi)^{c/2}\Gamma(2+c/2)D_r^{(n)}/(L^{\mathrm{rel}}_{3,n})^{c}=C^{(3)}_{n,r}$. Contrapositively, $\ell^f>C^{(3)}_{n,r}$ implies $\ell\nmid h_{3,n}/h_{3,n-1}$.

\emph{Remark (what changes relative to \cite{MO13}, \cite{MO16}).} The saturation input is that of \cite{MO13}; the floor is that of \cite{MO16} Lemma~2.5(2), applicable because $\operatorname{Nr}\eta_n=1$ transfers to the saturation root (Lemma~\ref{lem:normone}); the change is the convex body: \cite{MO13} \S4--5 and \cite{MO16} \S4 apply Minkowski's or Blichfeldt's theorem to an ideal lattice in coordinates and bound the height of $\delta_n^{\alpha}$ through the Mahler measure, which produces the factors $c!$ resp.\ $((c+2)/2)!\,(\sqrt{2\pi}/(3^{3/4}\log(\cdots)))^{c}$ in $G_1(3,s,f)^f$ and $G(3,s,f)^f$; here Blichfeldt's theorem acts on the log lattice with its exact covolume $D_r^{(n)}$ (Theorem~\ref{thm:rank3}), and the floor enters through $(L^{\mathrm{rel}}_{3,n})^{c}$ with $L^{\mathrm{rel}}_{3,n}$ growing like $\sqrt{3^n}$.
